\begin{proposition}[Franel pullback and cover-adapted coordinates]
\label{prop:hasse-franel-pullback}
Let
\[
 \mathcal A(a)=\sum_{n\geq0}b_na^n,
 \qquad
 \mathcal F(x)=\sum_{n\geq0}f_nx^n,
 \qquad
 f_n=\sum_{k=0}^n\binom nk^3,
\]
and, for a prime $p$, let
\[
 A_p(a)=\sum_{m=0}^{p-1}b_ma^m,
 \qquad
 K_p(x)=\sum_{m=0}^{p-1}f_mx^m
 \quad\text{in }\mathbb F_p[a]\text{ and }\mathbb F_p[x],
\]
respectively.  Put
\[
 \phi(x)=\frac{x(1-8x)}{1+x}
 \quad\text{and}\quad
 \Psi_{p,m}(x)=x^m(1-8x)^m(1+x)^{p-1-m}.
\]
Then the following identity holds in $\mathbb F_p[x]$:
\begin{equation}
 \boxed{K_p(x)^2=(1+x)^{p-1}A_p(\phi(x))
       =\sum_{m=0}^{p-1}b_m\Psi_{p,m}(x).}
 \label{eq:hasse-franel-pullback}
\end{equation}
The $p$ polynomials $\Psi_{p,0},\ldots,\Psi_{p,p-1}$ are linearly
independent.  Consequently, if $[P]_{\Psi_{p,m}}$ denotes the $m$th
coordinate of a polynomial $P$ in their span, then
\begin{equation}
 [K_p^2]_{\Psi_{p,m}}=b_m
 \quad(0\leq m<p),
 \qquad
 Z(p)=\#\{m:[K_p^2]_{\Psi_{p,m}}=0\}.
 \label{eq:hasse-franel-coordinates}
\end{equation}
In particular, this statement concerns coordinates in the
cover-adapted basis, not the monomial coefficients of~$K_p^2$.
\end{proposition}

\begin{proof}
The quadratic transformation used in
\cite[Section~2]{CFVZ2026} is the characteristic-zero formal identity
\begin{equation}
 \mathcal A(\phi(x))=(1+x)\mathcal F(x)^2.
 \label{eq:hasse-franel-characteristic-zero}
\end{equation}
Both coefficient sequences have the $p$-Lucas property.  Hence, in
$\mathbb F_p[[a]]$ and $\mathbb F_p[[x]]$,
\[
 \mathcal A(a)=A_p(a)\mathcal A(a)^p,
 \qquad
 \mathcal F(x)=K_p(x)\mathcal F(x)^p.
\]
Substitute $a=\phi(x)$ into the first congruence and use
\eqref{eq:hasse-franel-characteristic-zero}.  On the one hand,
\[
 (1+x)\mathcal F(x)^2
 =A_p(\phi(x))(1+x)^p\mathcal F(x)^{2p};
\]
on the other hand it equals
$(1+x)K_p(x)^2\mathcal F(x)^{2p}$.  Cancellation in the integral domain
$\mathbb F_p[[x]]$ proves the first equality in
\eqref{eq:hasse-franel-pullback}.  Clearing the powers of $1+x$ gives the
displayed polynomial expansion; both sides have degree at most $2p-2$.

Finally, $\Psi_{p,m}$ has $x$-adic order exactly $m$ and leading
coefficient~$1$.  The transition matrix from
$(\Psi_{p,0},\ldots,\Psi_{p,p-1})$ to the monomials is therefore triangular
with diagonal entries~$1$.  This proves linear independence and the
coordinate identities.
\end{proof}

\begin{proposition}[The two normalized Ap\'ery square roots]
\label{prop:hasse-apery-square-roots}
Let $\Delta(a)=a^2-34a+1$ and let $C_+(a)=\sum_{r\geq0}c_ra^r$ and
$C_-(a)=\sum_{r\geq0}d_ra^r$ be defined by $c_{-1}=d_{-1}=0$,
$c_0=d_0=1$, and
\begin{align}
 4(r+1)^2c_{r+1}
 &=2(68r^2+34r+5)c_r-(2r-1)^2c_{r-1},
 \label{eq:hasse-square-root-plus}\\
 4(r+1)^2d_{r+1}
 &=2(68r^2+102r+39)d_r-(2r+1)^2d_{r-1}.
 \label{eq:hasse-square-root-minus}
\end{align}
Thus $c_1=5/2$ and $d_1=39/2$.  In $\mathbb Q[[a]]$ one has
\begin{equation}
 \mathcal A(a)=C_+(a)^2=\Delta(a)C_-(a)^2.
 \label{eq:hasse-universal-square-roots}
\end{equation}

For every prime $p\geq5$, put
\[
 \varepsilon_p=\frac{1-\left(\frac{-6}{p}\right)}2.
\]
The factorization theorem of
Caruso--F\"urnsinn--Vargas-Montoya--Zudilin
\cite[Theorem~1.2]{CFVZ2026} gives a unique polynomial $B_p$ normalized by
$B_p(0)=1$ such that
\begin{equation}
 A_p(a)=\Delta(a)^{\varepsilon_p}B_p(a)^2.
 \label{eq:hasse-CFVZ-factorization}
\end{equation}
It is explicitly the indicated truncation of a universal series:
\begin{equation}
 B_p(a)=
 \begin{cases}
  \displaystyle\sum_{r=0}^{(p-1)/2}\overline{c_r}a^r,
   &\left(\frac{-6}{p}\right)=1,\medskip\\
  \displaystyle\sum_{r=0}^{(p-3)/2}\overline{d_r}a^r,
   &\left(\frac{-6}{p}\right)=-1.
 \end{cases}
 \label{eq:hasse-square-root-truncations}
\end{equation}
Here the bars denote reduction modulo~$p$; the displayed rational
coefficients have denominators prime to~$p$ in the stated ranges.
\end{proposition}

\begin{proof}
The recurrences are the coefficient forms of the two second-order
operators
\begin{align*}
 L_+&=4\vartheta^2
 -2a(68\vartheta^2+34\vartheta+5)+a^2(2\vartheta+1)^2,\\
 L_-&=4\vartheta^2
 -2a(68\vartheta^2+102\vartheta+39)+a^2(2\vartheta+3)^2,
 \qquad \vartheta=a\frac{d}{da}.
\end{align*}
A direct symmetric-square calculation shows that $C_+^2$ is annihilated by
the Ap\'ery operator
\[
 \vartheta^3-a(2\vartheta+1)(17\vartheta^2+17\vartheta+5)
 +a^2(\vartheta+1)^3.
\]
Its constant and linear coefficients are $1$ and~$5$, so uniqueness of the
formal solution gives $C_+^2=\mathcal A$.  Conjugating $L_+$ by
$\Delta^{-1/2}$ gives $L_-$; equivalently,
$C_-=C_+/\sqrt\Delta$, with constant term~$1$.  This proves
\eqref{eq:hasse-universal-square-roots} and
\eqref{eq:hasse-square-root-minus}.

It remains only to identify the normalized factor in
\eqref{eq:hasse-CFVZ-factorization}.  In the square branch, comparison of
the coefficients of degrees at most $(p-1)/2$ with
$\mathcal A=C_+^2$ determines the coefficients of $B_p$ successively,
because the square-root map with constant term~$1$ is triangular with
diagonal coefficient~$2$.  This gives the first line of
\eqref{eq:hasse-square-root-truncations}.  Apply the same argument to
$\mathcal A/\Delta=C_-^2$ in the defect branch.  All divisions occurring
in the recurrences have index strictly below~$p$, so reduction modulo~$p$
is legitimate.
\end{proof}

\begin{remark}[Exact algebraic reduction versus the GCD conjecture]
\label{rem:hasse-franel-logical-status}
Equations~\eqref{eq:hasse-franel-coordinates} show that the assertion
\[
 \#\{m:[K_p^2]_{\Psi_{p,m}}=0\}=O(1)
\]
is exactly Hypothesis~\ref{hyp:Z}.  It is not equivalent to the pointwise
Ap\'ery GCD conjecture.  Indeed, Remark~\ref{rem:pointwise-gap} shows that
even abstract zero sets of cardinality one can align adversarially across
different primes, whereas the full conjecture is equivalent to the
pointwise radical estimate in Remark~\ref{rem:rainbow-open}.  Conversely,
no converse implication is known: the pointwise radical formulation permits
the individual cardinalities~$Z(p)$ to grow provided their positions do not
align across primes.  Moreover, the Poisson model of
Remark~\ref{rem:Zfalse} predicts that Hypothesis~\ref{hyp:Z} is false.

The pullback nevertheless gives an exact, Ap\'ery-specific residual
coefficient problem.  General $\ell$-adic trace bounds, monodromy, and
Dwork unit-root congruences do not by themselves imply anti-concentration
of these defining-characteristic coordinates.  This is a limitation of the
currently available general theorems, not a claim that Cartier or Dwork
methods cannot contribute after additional Ap\'ery-specific input.  For
the average results in this paper, the natural single-prime target is the
weaker estimate $\sum_{p\leq X}Z(p)\ll\pi(X)$; the all-index conjecture also
requires the cross-prime dispersion described above.
\end{remark}
